\documentclass[UTF8,openany]{ctexbook}
\usepackage[a4paper,margin=2.5cm]{geometry}
\usepackage{amsmath,amssymb,mathtools}
\usepackage{amsthm}
\usepackage{booktabs,longtable,array}
\usepackage{enumitem}
\usepackage{xcolor}
\usepackage{tikz}
\usepackage{tocloft}
\usepackage{hyperref}
% 全笔记统一数学宏
% 设计原则：
% 1. 用简短命令减少重复输入和排版差异；
% 2. 所有宏使用 \providecommand，避免覆盖已有定义；
% 3. 在可能产生歧义的地方保留显式写法（如 \mathbb{R}^d）。

% ---- 常用数集 ----
\providecommand{\R}{\mathbb{R}}
\providecommand{\N}{\mathbb{N}}
\providecommand{\Z}{\mathbb{Z}}
\providecommand{\Q}{\mathbb{Q}}
\providecommand{\C}{\mathbb{C}}

% ---- 概率与期望 ----
\providecommand{\E}{\mathbb{E}}
% 概率测度全书统一用斜体 P（见 STYLE_GUIDE.md 第 6 节）；\Prob 保留为兼容别名
\providecommand{\Prob}{P}
\providecommand{\Var}{\operatorname{Var}}
\providecommand{\Cov}{\operatorname{Cov}}
\providecommand{\Corr}{\operatorname{Corr}}
\providecommand{\Bias}{\operatorname{Bias}}
\providecommand{\SE}{\operatorname{SE}}
\providecommand{\MSE}{\operatorname{MSE}}
\providecommand{\ind}{\mathbf{1}}          % 指示函数

% ---- 优化算子 ----
\providecommand{\argmax}{\operatorname*{arg\,max}}
\providecommand{\argmin}{\operatorname*{arg\,min}}

% ---- 线性代数 ----
\providecommand{\trans}{\mathsf{T}}        % 转置
\providecommand{\conj}{\mathsf{H}}         % 共轭转置
\providecommand{\trace}{\operatorname{tr}}
\providecommand{\diag}{\operatorname{diag}}
\providecommand{\rank}{\operatorname{rank}}
\providecommand{\nullsp}{\operatorname{Null}}
\providecommand{\rangesp}{\operatorname{Range}}

% ---- 范数、绝对值、内积、集合 ----
% 这些宏自动加 \left...\right，但在行内小公式中可能偏宽。
% 如果追求紧凑，可以改用 \norm*、\abs*（需要 mathtools 的 \DeclarePairedDelimiter）。
\providecommand{\norm}[1]{\left\lVert#1\right\rVert}
\providecommand{\abs}[1]{\left|#1\right|}
\providecommand{\ip}[2]{\left\langle#1,#2\right\rangle} % inner product
\providecommand{\set}[1]{\left\{#1\right\}}
\providecommand{\given}{\mid}

% ---- 常用算子 ----
\providecommand{\sign}{\operatorname{sign}}
\providecommand{\relu}{\operatorname{ReLU}}
\providecommand{\softmax}{\operatorname{softmax}}
\providecommand{\logsumexp}{\operatorname{LSE}}
\providecommand{\KL}{D_{\mathrm{KL}}}
\providecommand{\TV}{\operatorname{TV}}

% ---- 标注与强调 ----
\providecommand{\eqdef}{\stackrel{\text{\tiny def}}{=}} % 按定义相等

\hypersetup{
  colorlinks=true,
  linkcolor=black,
  urlcolor=blue,
  citecolor=blue,
  bookmarksdepth=1,
  psdextra=true,
  pdftitle={Beyond Formulas},
  pdfsubject={从公式到结构、判断与可信计算}
}
\AtBeginDocument{%
\renewcommand{\cftchapdotsep}{1}%
\renewcommand{\cftchapleader}{\cftdotfill{\cftchapdotsep}}%
\renewcommand{\cftchapafterpnum}{\cftparfillskip}%
}
\pdfstringdefDisableCommands{%
  \let\, \space
  \let\quad \space
  \let\qquad \space
}
\setlist{nosep}
\setlength{\emergencystretch}{3em}
\raggedbottom
\setcounter{secnumdepth}{0}
\setcounter{tocdepth}{0}
% 无编号定理类环境（与全书无编号风格一致；用法见 STYLE_GUIDE.md 第 4 节）
\theoremstyle{definition}
\newtheorem*{definition}{定义}
\newtheorem*{example}{例}
\newtheorem*{remark}{注记}
\theoremstyle{plain}
\newtheorem*{theorem}{定理}
\newtheorem*{proposition}{命题}
\newtheorem*{lemma}{引理}
\providecommand{\tightlist}{\setlength{\itemsep}{0pt}\setlength{\parskip}{0pt}}
\newcounter{none}
\title{Beyond Formulas\\[0.8em]\large 从公式到结构、判断与可信计算}
\author{}
\date{}
\begin{document}
\chapter*{数值微分与积分：同一份采样，为何难度相反}
\phantomsection\label{ch:064}
导数观察局部变化，会把高频噪声和舍入误差放大；积分汇总区间贡献，能平均局部扰动。两类方法都可从"用低次多项式替代函数"推导，但误差传播完全不同。

以中心差商近似 \(f'(x)\) 为例。取 \(f(x)=\sin x\)，\(x=\pi/4\)，真值为 \(\cos(\pi/4)\approx 0.70710678\)。用步长 \(h=10^{-3}\)，差分得 \(0.70710666\)，约 8 位正确。把 \(h\) 缩小到 \(10^{-8}\)，截断误差可降至 \(10^{-16}\) 量级——但浮点舍入把 \((f(x+h)-f(x-h))/(2h)\) 的分子推向了消去区，有效数字反而只剩约 8 位。步长不是越细越好。

本单元从步长的 U 形误差开始，再把梯形、Simpson、Gaussian 与自适应求积连成一条选择路径。

\begin{enumerate}
\tightlist
\item
  数值微分为何不能无限缩小步长
\item
  从插值多项式到梯形与 Simpson 求积
\item
  Gaussian 与自适应求积
\end{enumerate}

同一个 \(h\) 扮演两种角色：变小降低截断误差，却也可能放大舍入、噪声和采样盲区。公式的阶数只描述进入渐近区后的理想行为；实际选择还要看函数是否光滑、能否自由取样，以及粗细两套节点是否可能同时漏掉结构。

前一单元：\hyperref[ch:063]{``插值与逼近：通过数据点，不等于理解点之间''}。

\subsection{本章篇目}

以下篇目按概念依赖排列；若从具体问题进入，也可以先打开最接近当前困惑的一篇，再沿篇内链接回补。

\begin{itemize}
\tightlist
\item \hyperref[sec:08-Numerical-Analysis/04-Numerical-Differentiation-and-Integration/01]{``数值微分为何不能无限缩小步长''}；
\item \hyperref[sec:08-Numerical-Analysis/04-Numerical-Differentiation-and-Integration/02]{``从插值多项式到梯形与 Simpson 求积''}；
\item \hyperref[sec:08-Numerical-Analysis/04-Numerical-Differentiation-and-Integration/03]{``Gaussian 与自适应求积''}。
\end{itemize}

\section{数值微分为何不能无限缩小步长}
\label{sec:08-Numerical-Analysis/04-Numerical-Differentiation-and-Integration/01}

学过微积分的人看到导数定义

\[
f'(x)=\lim_{h\to0}\frac{f(x+h)-f(x)}{h},
\]

脑子里会冒出一个很自然的程序思路：取一个极小的 \(h\)，算一下差分，除以 \(h\)，完事。如果 \(10^{-6}\) 不够精确，换 \(10^{-12}\)。再不行就换 \(10^{-16}\)——双精度的极限。极限说的是 \(h\to0\)，那程序就应该尽量逼近零，不是吗？

不是。我们用具体的数字来试。取 \(f(x)=\sin x\)，求 \(x=1\) 处的导数。真值是 \(\cos 1\approx 0.54030230586\)。前向差分的结果：

\[
\begin{array}{c|c|c}
h & \text{近似值} & \text{误差} \\
\hline
10^{-1} & 0.49736 & 4.3\times10^{-2} \\
10^{-2} & 0.53608 & 4.2\times10^{-3} \\
10^{-4} & 0.54030 & 2.6\times10^{-6} \\
10^{-6} & 0.540302 & 7.3\times10^{-9} \\
10^{-8} & 0.5403023 & \mathbf{1.4\times10^{-9}} \\
10^{-10} & 0.54030 & 1.8\times10^{-6} \\
10^{-12} & 0.539 & 1.0\times10^{-3} \\
10^{-14} & 0.55 & 1.0\times10^{-2} \\
10^{-16} & 0 & 5.4\times10^{-1} \\
\end{array}
\]

误差先下降，到 \(h\approx 10^{-8}\) 附近触底，之后不仅回升，还加速恶化，最后直接给出零。这不是计算错误——如果你在双精度下计算 \(\sin(1+10^{-16})\)，\(1+10^{-16}\) 距离下一个可表示的浮点数还差约 50 倍，被舍入成了 \(1\)。\(\sin(1+10^{-16})=\sin(1)\)，分子为零，除以任何数都是零。极限定义假设每一步算术都精确，程序没有这个前提。

这个 U 形误差曲线是本节的核心研究对象。两股力量在反方向拉扯：步长大时截断误差主宰，步长小时舍入误差爆发。最优步长不在 \(h\to0\) 的极限处，而在两股力量平衡的谷底。

\subsection{Taylor 展开给出差分公式的截断误差}

你要逼近的是 \(f'(x)\)。把 \(f(x+h)\) 在 \(x\) 处 Taylor 展开：

\[
f(x+h)=f(x)+hf'(x)+\frac{h^2}{2}f''(\xi),
\]

其中 \(\xi\) 在 \(x\) 和 \(x+h\) 之间。移项并除以 \(h\)：

\[
D_h^+f(x)=\frac{f(x+h)-f(x)}{h}
=f'(x)+O(h).
\]

这是前向差分，一阶精度，因为扔掉的项是 \(h\) 的一阶量。每把 \(h\) 减半，截断误差大约减半。

中心差分用到了 \(x\) 两侧的信息。同时展开 \(f(x+h)\) 和 \(f(x-h)\)：

\[
\begin{aligned}
f(x+h)&=f(x)+hf'(x)+\frac{h^2}{2}f''(x)+\frac{h^3}{6}f'''(\xi_+),\\
f(x-h)&=f(x)-hf'(x)+\frac{h^2}{2}f''(x)-\frac{h^3}{6}f'''(\xi_-).
\end{aligned}
\]

两式相减，\(f(x)\) 和 \(f''(x)\) 项因对称性精确抵消（设 \(f\in C^3\)，两个余项的中值点可由介值定理合并为一个 \(\xi\)）：

\[
D_h^0f(x)
=\frac{f(x+h)-f(x-h)}{2h}
=f'(x)+\frac{h^2}{6}f'''(\xi).
\]

中心差分的截断误差从一阶跳到二阶，用了两次函数值。多出来的一次函数值买的不是"多一个信息点"，而是买了对称性。对称性让偶次幂项在相减中自动消掉，留下来的第一个未抵消项就是 \(h^2\) 阶。这是许多数值方法反复出现的主题——通过对称放置样本点，免费获取一个额外精度阶。

在区间边界无法访问 \(x-h\) 时必须用单边公式；若使用超过两个节点，可以通过插值多项式求导构造更高阶的差分模板（stencil）。更高阶的模板减少光滑函数的截断误差，但使用更宽的邻域和更大的正负交替权重——这意味着对噪声和边界效应更敏感。阶数不是免费的午餐。

\subsection{总误差的 U 形：两股力量的交汇}

截断误差是 \(h\) 减小时变好的一方。但浮点计算中还有另一方。每个函数值含绝对舍入误差约 \(C_ru\)，其中 \(u\) 是机器舍入单位——双精度下约 \(1.1\times10^{-16}\)，\(C_r\) 约等于函数值本身的量级。计算差分 \(f(x+h)-f(x)\) 时，两个接近数的相减让有效位数大量丢失。再除以 \(2h\)，残存的舍入误差被 \(1/h\) 放大。

以中心差分为例，总误差的粗略模型是

\[
E(h)\approx C_th^2+\frac{C_ru}{h}.
\]

第一项 \(C_th^2\) 随着 \(h\) 减小而减小——这就是截断误差在改善。第二项 \(C_ru/h\) 随着 \(h\) 减小而增大——舍入相消的代价被小分母放大。两者在相反的斜率上运动，交叉处就是总误差的最小值。

平衡两项：令 \(C_th^2=C_ru/h\)，解得 \(h^3=C_ru/C_t\)，给出主导量级

\[
h\asymp u^{1/3},
\]

其中 \(\asymp\) 表示量级相同，只保留了主导幂次。对于双精度，\(u^{1/3}\approx 4.8\times10^{-6}\)，这解释了为何中心差分的最优步长大约在 \(10^{-5}\) 到 \(10^{-6}\) 量级——不是 \(10^{-16}\)。

前向差分的截断误差是 \(O(h)\) 而非 \(O(h^2)\)，误差模型变为

\[
E(h)\approx C_th+\frac{C_ru}{h}.
\]

平衡得到 \(h\asymp u^{1/2}\)。双精度下约 \(10^{-8}\)——比中心差分的最优步长更小，但绝对误差却更大（因为截断误差随 \(h\) 只线性下降而非平方下降）。中心差分既更精确，又可以承受更大的实用步长——后者在函数的二阶导数取之不易时反而是优势。

这些指数不是无条件默认步长。如果你的函数值来自仿真容差或物理测量，那么有效的"噪声单位"不是机器精度 \(u\) 而是测量噪声的幅度 \(\eta\)——把 \(u\) 换成 \(\eta\)，最优 \(h\) 会显著右移。越粗糙的数据，越不能信任小步长。

\subsection{微分天然放大高频}

设你的观测包含快速振荡的小噪声：

\[
y(x)=f(x)+\varepsilon\sin(\omega x).
\]

噪声的幅度只有 \(\varepsilon\)，看似无害。但噪声的导数是 \(\varepsilon\omega\cos(\omega x)\)——幅度变成了 \(\varepsilon\omega\)。频率 \(\omega\) 越高，微分放大越强。一个幅度 \(10^{-3}\)、频率 \(1000\) 的噪声，在数值导数里会变成幅度 1 的振荡——把真实信号全部淹没。

对含噪数据硬套高阶差分公式，你可能得到一个公式上非常光滑、结果却非常离谱的"导数"。高阶差分假设函数在模板宽度内被低次多项式良好逼近；噪声打破了这一前提。

处理含噪数据的更可靠方向是：把"估计导数"视为一个带正则化的局部建模问题。先用局部多项式拟合、平滑样条或物理模型提取可辨识的趋势，再对趋势模型求导。平滑程度会引入偏差（你刻意忽略了一些高频变化），却抑制了不可恢复的高频噪声。偏差—方差的权衡在此以另一种形态出现：越平滑，噪声抑制越好，但被误归为噪声的真实局部变化也越多。

\subsection{complex-step：绕开实数相消的迂回策略}

如果被求导函数 \(f\) 可以解析延拓到复平面，并且你的全部代码路径都支持复数运算，有一个巧妙的迂回路线：

\[
f(x+ih)=f(x)+ihf'(x)-\frac{h^2}{2}f''(x)+O(h^3).
\]

取虚部并除以 \(h\)：

\[
f'(x)\approx\frac{\operatorname{Im}f(x+ih)}{h}.
\]

注意这个公式里没有两个接近实数相减的操作——\(f(x+ih)\) 和 \(f(x)\) 不在同一条实轴上竞争有效位。因此 \(h\) 可以取到 \(10^{-100}\) 甚至机器下溢的极限，截断误差被压到机器精度的平方级别，却不会被舍入相消的 U 形曲线干扰。

但代价同样明确。第一，函数必须可解析延拓——绝对值、条件分支、max/min、取整、逻辑判断、非解析特殊函数，任何一处不可解析都会让回答失去意义。第二，所有被 \(f\) 调用的子程序、外部库、乃至底层数学原语都必须支持复数运算并保持解析性。第三，复数路径的成本大约是实数路径的 2 到 8 倍——函数体越长，开销越大。这不是一个"对所有黑盒函数按下就行的通用技巧"。

另一种避开有限差分的路线是自动微分（automatic differentiation）。它通过对程序基本运算的链式法则传播，累积出精确的局部导数值——既没有截断误差，也没有步长选择的困扰。但它要求程序由可微基本操作构成，且仍然受原始程序本身的浮点舍入和不可微点（如分支条件跳变）影响。有限差分、复步微分和自动微分各自的适用边界，应作为工程选型时的一项清醒判断。

\section{从插值多项式到梯形与 Simpson 求积}
\label{sec:08-Numerical-Analysis/04-Numerical-Differentiation-and-Integration/02}

你需要算 \(\int_0^1 e^{-x^2}\,dx\)。不定积分没有初等形式——大一微积分课上老师就告诉过你"这个积不出来"。但数值上你只需要一个六位有效数字的答案，不需要原函数。

你面前只有函数在离散点上的值。你觉得应该用梯形法——把区间切成很多小段，每段上把曲线近似成直线，面积就是梯形。公式你在别处见过：\(h[f_0/2+f_1+\cdots+f_{n-1}+f_n/2]\)。但你又隐约记得还有个 Simpson 法，据说是把曲线近似成抛物线，精度高出不少。这两个公式看着完全不像——梯形权重是 \(1/2,1,1,\ldots,1,1/2\)，Simpson 权重是 \(1,4,2,4,2,\ldots,4,2,4,1\)——它们之间有什么联系吗？

本节要建立一个统一的视角：梯形法和 Simpson 法都在做同一件朴素的事——先用容易积分的低次多项式替代原函数，再把替代函数精确积掉。从这个视角看去，两者的权重来源、误差阶数、乃至更高阶方法的推广路径，全部清晰可见。

\subsection{求积公式的通用模板：插值，再积分}

要近似

\[
I=\int_a^bf(x)\,dx,
\]

在 \([a,b]\) 上选 \(n+1\) 个节点 \(x_0,\ldots,x_n\)，做插值多项式 \(p_n(x)=\sum_{i=0}^nf(x_i)\ell_i(x)\)。因为插值多项式容易积分——它就是多项式的线性组合——直接对两边积分：

\[
\int_a^bf(x)\,dx\approx\int_a^bp_n(x)\,dx
=\sum_{i=0}^nf(x_i)\underbrace{\int_a^b\ell_i(x)\,dx}_{w_i}.
\]

每个节点 \(x_i\) 获得一个权重 \(w_i\)，这个权重就是把对应的 Lagrange 基函数从 \(a\) 到 \(b\) 积分得到的一个数。权重只依赖节点位置和积分区间，与函数值 \(f(x_i)\) 无关。

这个模板的统一性值得强调：你不需要单独记忆梯形法是一种面积近似、Simpson 法是另一种——两者都是"选节点 \(\to\) 插值 \(\to\) 积分基函数 \(\to\) 得到权重"这个流水线的输出。不同之处仅在于用了几个节点、放在什么位置。把求积公式放回"先插值再积分"这个框架后，所有公式的系数都不再是来路不明的魔法数字。

\subsection{梯形法：用弦替代曲线，积分弦}

单区间 \([a,b]\) 的情形最简单。两个端点的线性插值就是连接 \((a,f(a))\) 和 \((b,f(b))\) 的弦。Lagrange 基函数：

\[
\ell_0(x)=\frac{b-x}{b-a},\qquad
\ell_1(x)=\frac{x-a}{b-a}.
\]

分别积分：

\[
w_0=\int_a^b\frac{b-x}{b-a}\,dx=\frac{b-a}{2},\qquad
w_1=\int_a^b\frac{x-a}{b-a}\,dx=\frac{b-a}{2}.
\]

于是

\[
T=\frac{b-a}{2}[f(a)+f(b)].
\]

就是梯形的面积公式——弦下方直角梯形的面积。

将 \([a,b]\) 均匀分成 \(n\) 段，步长 \(h=(b-a)/n\)，每个子区间独立使用梯形公式，相邻区间在共享端点处的权重相加。内点被两个相邻区间各贡献一次半权重权，首次出现时"借走"半份，下一次出现时"归还"半份，合起来正好是一份完整权重；首尾端点只属于一个区间，拿半份。复合梯形法：

\[
T_n=h\left[
\frac12f(x_0)+\sum_{i=1}^{n-1}f(x_i)+\frac12f(x_n)
\right].
\]

若 \(f\in C^2[a,b]\)，由每个子区间上线性插值误差的积分，存在 \(\xi\) 使

\[
I-T_n=-\frac{b-a}{12}h^2f''(\xi).
\]

全局二阶：网格点数加倍，误差约四分之一。对凸函数，每个子区间上弦都在函数图像上方（凸函数的图像弯在弦之下），梯形法高估。例如 \(f(x)=x^2\) 在 \([0,1]\) 上单梯形给出 \(\frac12\)，真值 \(\frac13\)。

非均匀节点不能套用统一的 \(h\) 权重，应逐区间用各自的宽度；对光滑周期函数且端点各阶导数匹配，复合梯形法的误差中所有带端点导数的项会反复抵消，实际精度可能远高于 \(O(h^2)\)——这个"梯形法的超收敛"只对周期光滑函数成立，不是普遍性质。

\subsection{Simpson 法：两个小区间共享一条抛物线，对称性白送一阶}

现在把眼光放宽到三个等距点：\(a\)、中点 \(m=(a+b)/2\)、\(b\)。在这三点上做二次 Lagrange 插值，积掉三条基函数：

\[
\ell_0(x)=\frac{(x-m)(x-b)}{(a-m)(a-b)},\quad
\ell_1(x)=\frac{(x-a)(x-b)}{(m-a)(m-b)},\quad
\ell_2(x)=\frac{(x-a)(x-m)}{(b-a)(b-m)}.
\]

换元 \(x=a+th\)，在 \(t\in[0,2]\) 上计算积分，得到权重：

\[
w_0=\frac{h}{3},\qquad w_1=\frac{4h}{3},\qquad w_2=\frac{h}{3}.
\]

因为 \(b-a=2h\)，公式常写为

\[
S=\frac{b-a}{6}
[f(a)+4f(m)+f(b)].
\]

中点权重是中点的 4 倍——它同时服务于左右两半段，抛物线的顶点信息被加倍征用。

均匀网格上复合 Simpson 法要求子区间数 \(n\) 为偶数（每两个相邻子区间共享一组三个节点做一条抛物线）：

\[
S_n=
\frac h3\left[
f_0+f_n
+4\sum_{i\text{ odd}}f_i
+2\sum_{i\text{ even},\,0<i<n}f_i
\right].
\]

奇下标权重 \(4\)，内部偶下标权重 \(2\)——每对相邻子区间中，中点属于本抛物线的内部节点，独占对应基函数的全部积分；共享端点则是左右两条抛物线各贡献一半，合起来权重 \(2\)。

虽然由二次插值推导，Simpson 规则因对称性能精确积分三次多项式。理由如下：奇函数（如 \((x-m)^3\)）在 \([a,b]\) 上关于中点 \(m\) 积分为零；Simpson 的三点对称放置让三次奇函数对应的插值多项式精确捕捉到这一对称性。误差中 \(f'''\) 的贡献因此消失，第一个残留项来自 \(f^{(4)}\)。若 \(f^{(4)}\) 连续，

\[
I-S_n=-\frac{b-a}{180}h^4f^{(4)}(\xi),
\]

所以网格点数加倍，渐近误差约缩小 \(1/16\)。这和梯形法每加倍只缩小到 \(1/4\) 相比，是质的飞跃——同样的计算代价买到两个额外精度阶。

不规则网格上 \(1,4,2,\ldots,4,1\) 的权重模式不能照搬——等距是抛物线积分系数的前提。函数有尖点、跳跃或端点奇性时，高阶导数的存在性被破坏，四阶的误差结论退化回更低的实际阶数。

\subsection{Richardson 外推：梯形和 Simpson 本是一条渐近链上的相邻节点}

考虑梯形法用步长 \(h\) 和半步长 \(h/2\) 计算两次。如果梯形误差有渐近展开

\[
T_h=I+Ch^2+O(h^4),
\]

那么半步长为

\[
T_{h/2}=I+C\frac{h^2}{4}+O(h^4).
\]

现在你手里有两个近似值，已知其主导误差项之间恰有 4 倍关系。把第二个式子乘以 4，减第一个：

\[
\frac{4T_{h/2}-T_h}{3}
=I+\frac{4O(h^4)-O(h^4)}{3}
=I+O(h^4).
\]

\(h^2\) 阶的主误差项被消灭，剩余误差从二阶直接跳到四阶。这个组合恰好就是复合 Simpson 公式——不是近似意义上的"差不多"，而是代数上严格相等。

验证一下权重算术。记细网格步长 \(h'=h/2\)。粗网梯形 \(T_h\) 涉及点 \(x_0,x_2,x_4,\ldots\)（每隔一个点），细网梯形 \(T_{h/2}\) 涉及全部点 \(x_0,x_1,x_2,\ldots\)。在 \(\frac{4T_{h/2}-T_h}{3}\) 中：

\begin{itemize}
\tightlist
\item
  端点（如 \(x_0\)）：细网梯形贡献 \(4\cdot\frac{h'}{2}=2h'\)，粗网梯形贡献 \(-\frac{1}{3}\cdot h=-\frac{2h'}{3}\)，合为 \(\frac{4h'}{3}=\frac{h}{3}\cdot\frac{1}{2}\)，对应 Simpson 的端点权重 1；
\item
  奇下标内部点（只出现在细网）：\(4\cdot h'=\frac{4h}{3}\cdot\frac{1}{2}\)，对应 Simpson 权重 4；
\item
  偶下标内部点（两网共享）：细网 \(4\cdot\frac{h'}{2}=2h'\)，粗网 \(-\frac{1}{3}\cdot h=-\frac{2h'}{3}\)，合为 \(\frac{4h'}{3}\)，再除以 \(\frac{h}{3}\) 对应 Simpson 权重 2。
\end{itemize}

全部权重与前面推导的 \(S_n\) 完全一致。Richardson 外推之所以有效，是因为它利用了你对误差幂次的知识——已知主导误差正比于 \(h^2\)，就用两个分辨率的组合将其消去。同样的逻辑可以继续向上：消去 \(h^4\) 项得到更高阶的求积公式（Romberg 积分），每一步都在消耗一个渐近幂次。

用 \(n\) 和 \(2n\) 的结果比较来估计误差是同样的原理，但前提是已经进入预期的渐近区。如果两次近似共同遗漏一个窄峰——比如函数在 \(x=0.713\) 处有一个宽度远小于 \(h\) 的尖峰——\(T_n\) 和 \(T_{2n}\) 可能数值非常接近，却同时错得离谱。外推帮不了你看见从未被采样到的结构。

\subsection{积分比微分温和，但温和有温和的陷阱}

对比上一节数值微分的 U 形误差曲线，积分的处境要友善得多。局部函数值误差在积分中乘上区间宽度后求和——宽度 \(h\) 本身就是一个小数，把误差压缩了一阶。微分则相反：相减后还要除以小 \(h\)，把误差膨胀一阶。从信号处理的角度，积分是平滑操作（把高频幅度除以频率），微分是放大高频的操作（把高频幅度乘以频率）。

因此"积分比微分温和"是一个可以定量信赖的陈述。但温和不意味着万能：

\begin{itemize}
\tightlist
\item
  端点奇性。如果被积函数在积分限处发散——例如 \(\int_0^1 x^{-1/2}dx\) 在 \(x=0\) 处——等距节点的多项式插值在奇点附近完全失效，需要专门处理（如变量替换、奇异减法、或 Gauss--Jacobi 求积）。
\item
  强振荡。\(\int_0^{2\pi}\sin(100x)f(x)dx\) 型的积分，正负值快速交替，等距采样可能需要远超直觉的节点数才能分辨。Filon 方法和 Levin 型方法为这类问题专门设计。
\item
  极窄峰。带宽远小于网格间距的尖峰可能完全漏过所有采样点，得到精确而完全错误的零。自适应求积通过局部误差估计检测到"这里还需要更细"。
\item
  无限区间。\(\int_{-\infty}^\infty\) 需要截断、变量压缩（如 \(x=\tan t\)）或 Gauss--Hermite 型求积。
\item
  高维积分。一维的 \(n\) 点求积公式推广到 \(d\) 维变成 \(n^d\) 点——维数灾难直接作用在求积节点数上。稀疏网格、Monte Carlo 和拟 Monte Carlo 方法为高维而存在。
\end{itemize}

梯形法和 Simpson 法是你工具箱里最常被调用的两把通用工具。记住它们的来源——插值再积分——你就能随时推导出非均匀网格上的对应公式、判断何时精度阶会退化、以及决定是否值得向 Romberg 外推或自适应求积迈出下一步。

\section{Gaussian 与自适应求积}
\label{sec:08-Numerical-Analysis/04-Numerical-Differentiation-and-Integration/03}

被积函数不总是一行公式。它可能是一次流体仿真的输出，一次昂贵的前向模拟，一次要排队等三十毫秒的远程调用。你的外层循环要算几十万个一维积分，摊到每个积分上的预算是三次函数求值。

上一节的工具在这个预算下能做到什么？三个等距点给出 Simpson 公式，对三次以内的多项式精确。想再准一点，只能加点——可是加点正是你付不起的。

问题的提法里其实藏着一个没被动过的自由度。Newton--Cotes 那一族公式先把节点钉在等距位置上，然后解出权重；权重是 \(n+1\) 个待定数，节点却被当成了常量。若把节点也放开，三次求值对应的是六个自由参数而不是三个。

\subsection{节点位置也是可以调的变量}

放开之后能换来多少，可以先看结果。

\begin{center}
\begin{tikzpicture}[x=1cm,y=1cm,font=\small]
  \node[right,font=\footnotesize,black!70] at (-0.1,4.15) {同样的函数求值次数，节点摆在哪里};
  \draw[black!45] (0,3.5) -- (5.4,3.5);
  \draw[black!45] (0,3.35) -- (0,3.65);
  \draw[black!45] (5.4,3.35) -- (5.4,3.65);
  \foreach \X in {0,2.7,5.4} \fill[black!80] (\X,3.5) circle (2.2pt);
  \node[right,font=\footnotesize,black!80] at (5.6,3.5) {等距三点};
  \draw[black!45] (0,2.85) -- (5.4,2.85);
  \draw[black!45] (0,2.7) -- (0,3.0);
  \draw[black!45] (5.4,2.7) -- (5.4,3.0);
  \foreach \X in {0.609,2.7,4.791} \fill[blue!75] (\X,2.85) circle (2.2pt);
  \node[right,font=\footnotesize,blue!75] at (5.6,2.85) {Gauss 三点};
  \draw[black!45] (0,2.2) -- (5.4,2.2);
  \draw[black!45] (0,2.05) -- (0,2.35);
  \draw[black!45] (5.4,2.05) -- (5.4,2.35);
  \foreach \X in {0.253,1.246,2.7,4.154,5.147} \fill[orange!85] (\X,2.2) circle (2.2pt);
  \node[right,font=\footnotesize,orange!90] at (5.6,2.2) {Gauss 五点};
  \node[right,font=\footnotesize,black!70] at (-0.1,1.6) {能被精确积掉的最高多项式次数};
  \draw[fill=black!30,draw=black!45] (0,1.02) rectangle (1.5,1.28);
  \node[right,font=\footnotesize,black!80] at (1.65,1.15) {\(3\)};
  \draw[fill=blue!30,draw=black!45] (0,0.57) rectangle (2.5,0.83);
  \node[right,font=\footnotesize,blue!75] at (2.65,0.7) {\(5\)};
  \draw[fill=orange!35,draw=black!45] (0,0.12) rectangle (4.5,0.38);
  \node[right,font=\footnotesize,orange!90] at (4.65,0.25) {\(9\)};
\end{tikzpicture}

\emph{上半部分是三套节点在同一个区间上的位置，下半部分是各自的代数精确度。三个点不动、只把位置从等距挪到 Gauss 点，精确度从三次升到五次；五个 Gauss 点到九次，而五个等距点（Boole 公式）仍旧只有五次。注意 Gauss 节点全部落在区间内部，且三点与五点两套除中心外没有一个重合。}
\end{center}

对一个正的权函数 \(w\)，Gauss 求积写成

\[
\int_a^bw(x)f(x)\,dx\approx\sum_{i=1}^n\omega_if(x_i),
\]

节点取为关于 \(w\) 的 \(n\) 次正交多项式的根，权重全为正。\(n\) 个点对次数不超过 \(2n-1\) 的多项式精确——这是任何 \(n\) 点公式所能达到的上限。图里那三根长短不一的条，量的就是这件事。

\subsection{正交多项式的根是那组最优落点}

上限先说清楚。任何一组给定的 \(n\) 个节点，取

\[
g(x)=\prod_{i=1}^n(x-x_i)^2,
\]

它是 \(2n\) 次多项式，在所有节点上取零，求积和必为零；而 \(g\ge0\) 且不恒为零，对正权函数它的真实积分为正。所以 \(2n\) 次是过不去的。

达到 \(2n-1\) 的构造只有一行。记 \(\pi_n\) 为关于 \(w\) 的 \(n\) 次正交多项式，任意 \(p\in P_{2n-1}\) 做带余除法：

\[
p=q\pi_n+r,\qquad q,r\in P_{n-1}.
\]

积分那边，\(q\pi_n\) 的积分因正交性为零（\(q\) 的次数低于 \(n\)）；求积和那边，节点取在 \(\pi_n\) 的根上，\(\pi_n(x_i)=0\)，这部分也为零。两边只剩 \(r\)，而 \(r\) 的次数不超过 \(n-1\)，让权重对 \(P_{n-1}\) 精确即可——那正是插值型求积公式的定义。

条件对账：正交性防的是积分那一侧留下 \(q\pi_n\) 的残项，要求 \(q\) 的次数严格低于 \(n\)，所以 \(p\) 的次数不能超过 \(2n-1\)；根落在开区间内部（正交多项式的根都是实的、单的、在支撑内）防的是节点跑到积分区间之外；权函数为正防的是权重出现负号——负权重会在求和中制造相消，把上一单元里那种有效位损失重新引进来。

区间 \([-1,1]\)、\(w=1\) 时正交多项式是 Legendre 多项式。\(P_2(x)=\tfrac12(3x^2-1)\) 的根给出两点公式 \(x_{1,2}=\pm1/\sqrt3\)，权重都是 \(1\)：

\[
\int_{-1}^1f(x)\,dx\approx f(-1/\sqrt3)+f(1/\sqrt3),
\]

两次求值精确积掉全部三次多项式。搬到一般区间 \([a,b]\) 要做仿射变换，节点按 \(x\mapsto\frac{a+b}{2}+\frac{b-a}{2}x\) 平移伸缩，权重同时乘 \(\frac{b-a}{2}\)；只换节点不换权重会让积分整体差一个尺度因子，是这类代码最常见的一处错。

权函数是可以挑的。把 \(e^{-x^2}\) 挑进 \(w\) 得到 Gauss--Hermite，把 \(e^{-x}\) 挑进去得到 Gauss--Laguerre，把端点奇性 \((1-x)^\alpha(1+x)^\beta\) 挑进去得到 Gauss--Jacobi。规则与测度对上号时，被积函数里最难缠的那部分被吸收进权重表，剩下的部分光滑，少数几个节点就够。计算高斯型期望时这一点特别值钱：\(\E[g(Z)]\) 对 \(Z\sim\mathcal N(0,1)\)，用五个 Hermite 节点常常比一万个随机样本还准。

Gauss 节点向区间两端加密的形态，与\hyperref[sec:08-Numerical-Analysis/03-Interpolation-and-Approximation/02]{``Runge 现象与节点选择''}里的 Chebyshev 节点是同一副长相。两处的理由也相通：等距布点在端点附近让节点乘积鼓出巨包，把节点往两端挪能把这个包压平。求积和插值在这件事上共用一个答案。

\subsection{代价是这组节点不能复用}

精度不是白拿的。Gauss 节点是无理数，位置随 \(n\) 整个改变——图里三点与五点两排，除了中心那个点，没有一个能对上。这意味着算完三点想加密到五点，前三次求值全部作废，要重头再算五次。等距的复合公式没有这个问题：从 \(n\) 段加密到 \(2n\) 段，原有的函数值一个不丢，只需补上新增的中点，这也正是上一节 Richardson 外推能成立的前提。

Gauss--Kronrod 就是为这笔账设计的补丁。它在 \(n\) 个 Gauss 节点之外再嵌入 \(n+1\) 个新节点，凑成一个 \(2n+1\) 点的规则；两套规则共享那 \(n\) 个已有的函数值，差值直接当误差指标用。代价是嵌入规则的代数精确度低于同点数的纯 Gauss 规则——拿一部分精度换一个免费的误差估计。几乎所有工业级的一维积分例程内部跑的都是它。

\subsection{自适应把预算搬到函数难缠的那一段}

前面讨论的都是固定网格：节点在哪由公式决定，与被积函数长什么样无关。可被积函数的困难往往集中在很小的一块。

\begin{center}
\begin{tikzpicture}[x=1.7cm,y=1.5cm,font=\small,>=stealth]
  \draw[black!30] (-0.05,0) -- (4.2,0);
  \draw[black!45] (0,0) -- (0,1.55);
  \draw[black!85,very thick] plot[smooth] coordinates {
    (0,0.250) (0.8,0.250) (1.5,0.250) (2.0,0.250) (2.1,0.255) (2.2,0.283)
    (2.3,0.406) (2.4,0.723) (2.5,1.171) (2.6,1.400) (2.7,1.171) (2.8,0.723)
    (2.9,0.406) (3.0,0.283) (3.1,0.255) (3.2,0.250) (3.6,0.250) (4.0,0.250)};
  \draw[black!35,dashed] (2.2,0) -- (2.2,1.45);
  \draw[black!35,dashed] (3.0,0) -- (3.0,1.45);
  \foreach \t in {0,0.25,...,4} \fill[black!70] (\t,-0.30) circle (1.4pt);
  \foreach \t in {0,0.6,1.2,1.8,2.1,2.2,2.3,2.375,2.45,2.525,2.6,2.675,2.75,2.825,2.9,3,3.2,3.6,4}
    \fill[blue!75] (\t,-0.62) circle (1.4pt);
  \node[right,font=\footnotesize,black!75] at (4.15,-0.30) {均匀 \(17\) 点};
  \node[right,font=\footnotesize,blue!75] at (4.15,-0.62) {自适应 \(19\) 点};
  \node[font=\footnotesize,black!70] at (1.05,1.15) {平坦段};
  \node[font=\footnotesize,black!70] at (3.75,1.15) {平坦段};
\end{tikzpicture}

\emph{几乎相同的求值预算，两种花法。均匀网格把十三个点撒在两侧的平坦段上，落进虚线之间那个窄峰的只有四个；自适应网格在平坦段几点带过，把十个点压进峰区。峰的面积占全积分的大半，误差也几乎全在那里。}
\end{center}

自适应求积的机制很简单：在每个区间上算一粗一细两个近似，用两者之差估计局部误差；达标就收下，不达标就把区间劈成两半，对每一半重复。图里那种点的分布不是谁摆出来的，是这个递归自己长出来的。

误差估计可以借上一节的渐近关系。设整个区间的一次 Simpson 结果是 \(S\)，左右半区间之和是 \(S_L+S_R\)。四阶意味着步长减半误差降到十六分之一，于是 \(I-S\approx16(I-S_L-S_R)\)，解出

\[
I-(S_L+S_R)\approx\frac{(S_L+S_R)-S}{15}.
\]

右边全是已经算出来的量。这个 \(1/15\) 就是 Richardson 外推倒过来用：那边拿幂次关系去消掉误差，这边拿同一个幂次关系去读出误差。Gauss--Kronrod 的误差指标是另一种读法，用两套阶数不同的规则代替两种步长。

实现时还有一串必须显式设定的东西：绝对容差与相对容差各一个（只设相对容差会在积分值接近零时永不收敛）、最大求值次数、递归深度上限、最小区间宽度，以及 \(\infty\) 与 NaN 出现时的处置。这些不是工程琐碎，是防住递归在奇点处无限细分下去的唯一手段。同一套自适应逻辑在\hyperref[sec:08-Numerical-Analysis/06-ODE/02]{``Runge--Kutta 与自适应步长''}里控制的是时间步长，控制器的骨架完全一样。

误差估计本身也会骗人。粗细两个结果接近，只说明加密这一次没带来变化，不说明结果对。若两套节点都从一个极窄的峰旁边跨过去，差值会很小，算法安心地收下一个错误答案——这与上一节比较 \(T_n\) 与 \(T_{2n}\) 时的陷阱是同一个。防它的办法不在数值方法内部：知道峰在哪就手动拆区间，或者先粗扫一遍定位可疑区域。

\subsection{不光滑时高阶反而吃亏}

代数精确度描述的是一个测试空间：规则在多项式上表现如何。被积函数离多项式很远时，这个指标失去参考价值。

Gauss 规则的误差里带着 \(f^{(2n)}\)。函数光滑、高阶导数有界时，这一项随 \(n\) 迅速衰减，收敛快得惊人；解析函数上甚至是几何速率，这条线索在\hyperref[sec:08-Numerical-Analysis/08-Fourier-and-Spectral-Methods/04]{``谱方法：用光滑性换高精度''}里还会展开。可一旦函数在区间内有跳跃、有尖点、有端点奇性，\(f^{(2n)}\) 根本不存在，高阶规则不再有任何优势，甚至比低阶更糟——它在更宽的邻域内假定一个高次多项式能代表函数，而这个假定错得更远。

这种时候正确的动作是降阶加密，不是升阶。分段低次公式加上足够多的段，收敛速度由函数实际的光滑度决定，不会因为方法的野心而恶化。更好的做法是先把不光滑处摘出来：在跳跃点把区间拆开，让每一段内部光滑；对已知的端点奇性做变量替换或奇性减除；把奇性吸进权函数，交给匹配的 Gauss--Jacobi 规则。

选规则的时候，这几条大致够用。手上只有一张离散数据表、无法在任意点求值，那么分段梯形是最诚实的选择——你没有凭空使用未采样曲率的权利。函数光滑且可以随意取样，Gauss 或 Clenshaw--Curtis 用很少的点就能到高精度。困难集中在局部，交给自适应。被积函数是解析的周期函数且积分覆盖整周期，复合梯形法反而异常优秀，因为误差展开中所有带端点导数的项两两抵销。端点奇性或无限区间，先做变换或换匹配的权函数，再谈阶数。

而一维的所有讨论到高维都要重来。张量积把 \(n\) 点公式变成 \(n^d\) 点，维数灾难直接落在求值次数上，稀疏网格能缓解一部分，超过十几维就得换到 Monte Carlo 那一类不依赖光滑度、收敛率与维数无关的方法上去。求积规则的精妙全建立在``函数在低维网格上可被多项式代表''这个前提上，前提不成立时，精妙本身就是负担。

\end{document}
